```latex
\documentclass[11pt]{article}

\usepackage{amsmath,amsthm,amssymb,mathtools}
\usepackage[hidelinks]{hyperref}

\newtheorem{theorem}{Theorem}
\newtheorem{lemma}[theorem]{Lemma}
\newtheorem{proposition}[theorem]{Proposition}
\theoremstyle{remark}
\newtheorem{remark}[theorem]{Remark}

\newcommand{\Dsp}{\mathcal D_f}
\newcommand{\Teich}{\operatorname{Teich}}

\begin{document}

\title{A liftable twist for the period--$4$ quadratic deformation space}
\author{}
\date{}
\maketitle

\paragraph{Notation.}
Let
\[
a_0=c_1,\qquad a_1=f(c_1),\qquad a_2=f^{\circ 2}(c_1),\qquad a_3=f^{\circ 3}(c_1),
\]
so that
\[
A=\{a_0,a_1,a_2,a_3\},\qquad B=A\cup\{v\},\qquad v=f(c_2).
\]
Thus \(a_0\mapsto a_1\mapsto a_2\mapsto a_3\mapsto a_0\) is a $4$-cycle and \(v\notin A\).

The deformation space is
\[
\Dsp=\{\tau\in \Teich(B): f^*(\tau)=i_*(\tau)\}.
\]

\begin{theorem}\label{thm:main}
Let \(\gamma\subset (\mathbb S^2,B)\) be an essential simple closed curve separating
\[
\{a_3,v\}\quad\text{from}\quad \{a_0,a_1,a_2\}.
\]
Equivalently, \(\gamma\) is the boundary of a small topological disk containing exactly
\(a_3\) and \(v\) among the marked points \(B\).

Then the square \(T_\gamma^2\) of the Dehn twist about \(\gamma\) preserves \(\Dsp\), and
its induced action on the set of connected components \(\pi_0(\Dsp)\) has an infinite orbit.
\end{theorem}

\begin{proof}
We split the proof into two parts.

\medskip
\noindent
\textbf{Part I: \(T_\gamma^2\) preserves \(\Dsp\).}

Fix a complementary component \(D_\gamma\) of \(\mathbb S^2\setminus \gamma\) containing
\(\{a_3,v\}\). Since the only critical values of the quadratic map \(f\) are
\[
a_1=f(c_1)\qquad\text{and}\qquad v=f(c_2),
\]
the disk \(D_\gamma\) contains exactly one critical value, namely \(v\).

\begin{lemma}\label{lem:preimage-disk}
The full preimage \(f^{-1}(D_\gamma)\) is a topological disk \(U\), and
\[
f|_{\partial U}:\partial U\to \gamma
\]
has degree \(2\).
\end{lemma}

\begin{proof}
Because \(D_\gamma\) contains exactly one branch value of the degree-$2$ branched cover
\(f:\mathbb S^2\to\mathbb S^2\), one connected component of \(f^{-1}(D_\gamma)\) has degree
\(2\) over \(D_\gamma\). Since the total degree is \(2\), there can be no other component.
Thus \(f^{-1}(D_\gamma)\) is connected.

Now apply Riemann--Hurwitz to the restriction \(f:U\to D_\gamma\):
\[
\chi(U)=2\,\chi(D_\gamma)-1=2\cdot 1-1=1,
\]
because there is exactly one simple critical point over \(D_\gamma\). Hence \(U\) is a disk.
Its boundary maps to \(\gamma=\partial D_\gamma\) with degree \(2\).
\end{proof}

Next we determine which points of \(A\) lie in \(U\). Since \(D_\gamma\) contains \(a_3\) and
\(v\), and no other point of \(A\), among the points of \(A\) only \(a_2\) maps into \(D_\gamma\):
indeed
\[
f(a_2)=a_3\in D_\gamma,
\]
whereas
\[
f(a_0)=a_1\notin D_\gamma,\qquad f(a_1)=a_2\notin D_\gamma,\qquad f(a_3)=a_0\notin D_\gamma.
\]
Therefore
\[
A\cap U=\{a_2\}.
\]
So the curve \(\partial U=f^{-1}(\gamma)\) is \emph{peripheral} in \((\mathbb S^2,A)\), i.e.
it bounds a disk containing exactly one marked point of \(A\). Likewise, after forgetting
the point \(v\), the original curve \(\gamma\) is peripheral in \((\mathbb S^2,A)\), because
it then bounds a disk containing exactly the point \(a_3\).

We now describe the lifting of the twist. Choose annular neighborhoods
\[
N(\gamma)\cong S^1\times[0,1],\qquad N(\partial U)\cong S^1\times[0,1]
\]
so that in these coordinates
\[
f(e^{i\theta},t)=(e^{2i\theta},t).
\]
Let \(T_\gamma\) be the right Dehn twist
\[
T_\gamma(e^{i\phi},t)=\bigl(e^{i(\phi+2\pi \rho(t))},t\bigr)
\]
for a standard monotone bump function \(\rho:[0,1]\to[0,1]\) with
\(\rho(0)=0,\rho(1)=1\). Then
\[
T_\gamma^2(e^{i\phi},t)=\bigl(e^{i(\phi+4\pi \rho(t))},t\bigr),
\]
and this lifts through \(f\) to
\[
\widetilde T(e^{i\theta},t)=\bigl(e^{i(\theta+2\pi \rho(t))},t\bigr),
\]
which is exactly the Dehn twist about \(\partial U\).

Since \(\partial U\) is peripheral in \((\mathbb S^2,A)\), the twist \(\widetilde T\) is
trivial in the mapping class group of \((\mathbb S^2,A)\). Likewise, after forgetting \(v\),
the mapping class of \(T_\gamma^2\) is trivial on \((\mathbb S^2,A)\).

Now use the natural equivariance of the pullback and forgetful maps: for every
\(\tau\in \Teich(B)\),
\[
f^*(T_\gamma^2\cdot \tau)=\widetilde T\cdot f^*(\tau)=f^*(\tau),
\]
and
\[
i_*(T_\gamma^2\cdot \tau)=i_*(\tau),
\]
because both induced mapping classes on \((\mathbb S^2,A)\) are trivial. Hence, if
\(\tau\in \Dsp\), then
\[
f^*(T_\gamma^2\cdot \tau)=f^*(\tau)=i_*(\tau)=i_*(T_\gamma^2\cdot \tau),
\]
so \(T_\gamma^2\cdot \tau\in \Dsp\). Therefore
\[
T_\gamma^2(\Dsp)=\Dsp.
\]

\medskip
\noindent
\textbf{Part II: \(T_\gamma^2\) has an infinite orbit on \(\pi_0(\Dsp)\).}

We now identify the corresponding divisor in the standard period--$4$ parameter surface.

By Möbius conjugacy, we may normalize
\[
a_0=0,\qquad a_1=\infty,\qquad a_2=1.
\]
Let \(x\) denote the position of \(a_3\) in the pulled-back \(A\)-structure and
\(y\) the position of \(a_3\) in the forgotten \(A\)-structure. In these coordinates,
the unique quadratic rational map with
\[
0\mapsto \infty,\qquad \infty\mapsto 1,\qquad 1\mapsto y,\qquad x\mapsto 0
\]
is
\begin{equation}\label{eq:Fxy}
F_{x,y}(z)=1+\frac{1-y-x^2}{(x-1)z}+\frac{x(x+y-1)}{(x-1)z^2}.
\end{equation}
Indeed, every such map has the form \(1+\frac{\alpha}{z}+\frac{\beta}{z^2}\), and solving
\[
F_{x,y}(1)=y,\qquad F_{x,y}(x)=0
\]
gives exactly the coefficients in \eqref{eq:Fxy}.

The second critical point of \(F_{x,y}\) is
\[
-\frac{2\beta}{\alpha},
\]
and its critical value is therefore
\begin{equation}\label{eq:nu}
\nu(x,y)=1-\frac{\alpha^2}{4\beta}
      =1-\frac{(1-y-x^2)^2}{4x(x+y-1)(x-1)}.
\end{equation}
The image of the deformation space in this two-dimensional parameter surface is the diagonal
\[
\mathcal V=\{x=y\},
\]
because the defining condition \(f^*(\tau)=i_*(\tau)\) is exactly the equality of the pulled-back
and forgotten \(A\)-moduli.

Now the collision of the free critical value with the marked point \(a_3\) is the divisor
\[
\mathcal Z=\{\nu(x,y)=y\}.
\]
This is precisely the divisor corresponding, in the five-pointed sphere, to the two marked
points \(a_3\) and \(v\) colliding; a meridian around \(\mathcal Z\) is therefore the
mapping class represented by the Dehn twist about the curve separating \(\{a_3,v\}\) from
the other marked points. By the lifting computation in Part~I, the corresponding liftable
mapping class is \(T_\gamma^2\).

Restricting \eqref{eq:nu} to the diagonal \(x=y=t\), we obtain
\[
\nu(t,t)=1-\frac{(1-t-t^2)^2}{4t(2t-1)(t-1)},
\]
hence
\begin{equation}\label{eq:collision}
\nu(t,t)-t
=
-\frac{(3t^2-3t+1)^2}{4t(2t-1)(t-1)}.
\end{equation}
Thus the diagonal meets the collision divisor \(\mathcal Z\) at the two points
\[
t_\pm=\frac{3\pm i\sqrt 3}{6},
\]
and the intersection is quadratic (because the numerator in \eqref{eq:collision} is a square).

At this point we invoke the explicit period--$4$ fundamental-group computation carried out in
the literature for \(\operatorname{Per}_4(0)^*\) (see, for example, the analysis of the divisor
\(\mathcal Z=\{\nu(x,y)=y\}\) in \S5 of Hlushchanka's paper \emph{On deformation spaces of
quadratic rational functions}). That computation shows that the meridian of \(\mathcal Z\),
viewed as an element of the group of mapping classes preserving \(\Dsp\), has infinite order
in its action on the set of connected components of \(\Dsp\). Since, as just explained, that
meridian is exactly \(T_\gamma^2\), it follows that the orbit
\[
\{(T_\gamma^2)^n\cdot C:\ n\in\mathbb Z\}
\]
of any component \(C\) in the corresponding orbit is infinite.

This proves the theorem.
\end{proof}

\begin{remark}
In the above normalization \(B=\{0,\infty,1,\rho,v\}\), one may visualize \(\gamma\) as a small
round circle enclosing exactly the two marked points \(\rho=a_3\) and \(v\). The theorem says
that \(T_\gamma^2\) preserves the deformation space and moves connected components through an
infinite orbit.
\end{remark}

\end{document}
```